\section{Introduzione}

Verranno trattati metodi e tecniche formali per specificare, disegnare e analizzare sistemi complessi. In particolare, \textbf{sistemi concorrenti e distribuiti}, costituiti da componenti che operano in modo indipendente e che interagiscono tra loro. Ci concentreremo sulle fasi di \textbf{specifica del problema} e \textbf{modellazione di una soluzione}. L'uso di metodi formali in queste fasi permette di verificare e validare il modello prima di procedere all'implementazione.
\\
\\
Per fare ciò verrano introdotti:
\begin{itemize}
	\item Un \textbf{linguaggio logico} per specificare le proprietà di comportamento di tali sistemi
	\item Le \textbf{reti di Petri} come strumento per modellare i sistemi concorrenti e analizzarne le proprietà
	\item Nozioni di base sui sistemi dinamici a tempi discreti (cenni ad automi cellulari)
\end{itemize}

\subsection{Teoria Generale delle Reti di Petri}

L'obiettivo di Petri era sviluppare una teoria matematica fondata sui principi della fisica moderna (teoria della relatività e fisica quantistica), che sia una \textbf{teoria dei sistemi} in grado di descrivere il \textbf{flusso di informazione} e permetta di analizzare \textbf{sistemi con organizzazione complessa}. Una teoria generale dei sistemi che processano informazione basata su: comunicazione, sincronizzazione, flusso di informazione, etc.
